\section{Methodology}\label{sec:method}

\subsection{Dataset Construction}\label{sec:dataset}

We extracted functions from ten open-source repositories at a particular commit (eight Java
projects: \texttt{commons-cli}, \texttt{commons-csv}, \texttt{commons-collections},
\texttt{commons-math}, \texttt{gson}, \texttt{jsoup}, \texttt{joda-time},
\texttt{jfreechart}; two Python projects: \texttt{flask}, \texttt{requests}), using Lizard
1.23.0 to compute cyclomatic complexity (CC). We select the medium complexity band
$5 \le CC \le 10$, as in our previous work.

While complexity is a precondition for being a challenging test target, a study sampling on
complexity confounds two phenomena: the quality of tests produced by the generator, and
whether the sampled target is reachable by the generator at all. We therefore applied five
additional criteria, each fixed before any test was generated, and each considered necessary
for an LLM, a search-based generator, or a human to reach the target through the public API.

\begin{description}
  \item[F1 --- Non-trivial body.] $nloc \ge 3$. This filters out stubs, abstract methods,
    and interface signatures that specify no branches.
  \item[F2 --- Public or exported.] Java: declared \texttt{public}. Python: not prefixed
    with an underscore. A \texttt{private} method cannot be invoked directly by an external
    test class; reflection can reach it, but then the compiler no longer checks any name in
    the test (Section~\ref{sec:tiers}), so compilation ceases to be evidence of anything.
  \item[F3 --- Unambiguous name.] Java: the method name occurs exactly once in its declaring
    class. Python: the function name occurs exactly once in its module. An overloaded name
    cannot appear in a prompt without additional disambiguating context, and the compiler
    rejects the resulting call as ambiguous.
  \item[F4 --- Verified declaring class.] The enclosing class is found with a brace-depth
    stack that \emph{pops} on block exit. Taking ``the nearest enclosing class declaration''
    without popping would attribute a method to whichever nested class was most recently
    declared in the file.
  \item[F5 --- Dynamically resolvable] (Python). The module is imported, the qualified name
    is walked with \texttt{getattr}, and the result is confirmed callable. Syntactic
    resolution is not sufficient evidence for run-time binding.
\end{description}

Table~\ref{tab:funnel} shows the complete funnel. We report every count-reducing stage
rather than only the final count, because the attrition is informative: of the 1{,}848
functions that passed the complexity filter, only 474 (26\,\%) were both public and
unambiguously named. Around three quarters of medium-complexity functions in these projects
are not valid test targets for an external generator.

\begin{table}[t]
\centering
\caption{Sampling funnel. Each row reports how many candidates survive the criterion. F4
corrects the declaring-class attribution that F3 depends on rather than removing candidates
in its own right, so it has no separate row.}
\label{tab:funnel}
\begin{tabular}{llr}
\toprule
Stage & Criterion & Remaining \\
\midrule
F0 & $5 \le CC \le 10$ (Lizard)          & 1{,}848 \\
F1 & non-trivial body                    & 1{,}848 \\
F3 & unambiguous name                    & 753 \\
F2 & public / exported                   & 474 \\
F5 & dynamically resolvable (Python)     & 64 \\
\midrule
   & final pool                          & Java 406, Python 64 \\
   & sampled                             & Java 60, Python 60 \\
\bottomrule
\end{tabular}
\end{table}

From the final pool we sample 60 Java and 60 Python functions, stratified across
repositories (Java: 3--9 functions per project; Python: 32 from \texttt{requests}, 28 from
\texttt{flask}). Complexity is skewed towards the lower end of the band ($CC{=}5$: 56
functions), and the median function contains 17 lines of code. Of the Python targets, 34 are
instance methods and 26 are module-level functions. Every one of the 120 functions was
checked for run-time accessibility before generating tests. One Java target was later found
to violate F2 --- our own parser had recorded a \texttt{protected} method as public --- and
is excluded, so Java results are reported over $n = 59$ and Python over $n = 60$
(Section~\ref{sec:threats}).

\subsection{Test Generation}\label{sec:generation}

All tests are generated by \texttt{gpt-4o-mini-2024-07-18} at
\texttt{temperature}\,$=$\,$0$, \texttt{top\_p}\,$=$\,$1.0$,
\texttt{max\_tokens}\,$=$\,$2048$, in a \textbf{single call per function} with no feedback
loop and no retrieval. We compare two \emph{one-shot} prompt conditions. Both contain
exactly one worked exemplar (a trivial \texttt{add} function with its test), so both are
one-shot in the usual sense, and both issue one call. They differ in exactly one respect:

\begin{description}
  \item[V1 --- baseline.] The prompt contains the target function name and source. This is
    the prompt used in our earlier study.
  \item[V2 --- target-specified.] Identical to V1, with one block inserted before the task
    description: the fully qualified target name, its signature, and declaring class, plus
    --- for instance methods --- a constructor call that has been \emph{executed} and
    verified to produce a valid receiver object.
\end{description}

Holding the worked exemplar fixed is crucial: an earlier variant of the enriched prompt
accidentally omitted the exemplar and thus differed from V1 in two respects at once (it
dropped the exemplar and it added the metadata), so no causal attribution was possible from
it. We report that variant only as a secondary observation and label it zero-shot, since it
is one.

\subsection{Measurement: Four Tiers}\label{sec:tiers}

A single ``valid / invalid'' flag conflates outcomes that differ in kind, so we report four
nested tiers. Each subsequent tier implies the one before it, and --- importantly --- the
first two can be satisfied by a test suite that does not test the target at all.

\begin{description}
  \item[T1 --- Compiles / collects.] The test file is syntactically and type-correct.
  \item[T2 --- At least one green test.] At least one test passes against the unmodified
    source.
  \item[T3 --- Reaches the target.] Branch coverage $> 0$ within the target's line range
    $[\textit{start},\textit{end}]$.
  \item[T4 --- Detects faults.] Mutation score $> 0$ within the same line range.
\end{description}

Only T4 speaks to a suite's ability to test the target; we demonstrate that T1 and T2 can be
satisfied by suites that verify nothing:

\begin{itemize}
  \item \textbf{T1 under reflection.} Reflection turns every name into a string, so the
    compiler stops checking it. A call to
    \texttt{getDeclaredMethod("aNameThatDoesNotExist")} compiles with return code $0$ and
    fails only at run time. Any compilation-based validity rate is therefore uninformative
    for suites that use reflection.
  \item \textbf{T2 under an assertion-free suite.} The search-based suite generated for
    \texttt{requests.hooks} (target \texttt{dispatch\_hook}) contains exactly one test, and
    that test is a single bare call to a \emph{different} function in the same module, with
    no assertion of any kind. It compiles, it collects, and it passes --- T1 and T2 are both
    satisfied --- while measured branch coverage inside the target's line range is $0$. A
    green-test criterion records this as a success.
\end{itemize}

The two tiers are therefore not merely weaker than T4; they can be satisfied by suites that
are, respectively, unchecked by the compiler and empty of assertions.

Branch coverage is measured with \texttt{coverage.py} (Python), attributed to the target's
line range rather than to the whole file; T3 is therefore reported for Python only. The Java
pipeline executes the generated JUnit suites through the JUnit Platform launcher and rests
on T4, so no Java branch-coverage figure is claimed anywhere in this paper. Mutation analysis
applies the same operator set to both languages --- arithmetic ($+/-$, $\times/\div$),
relational ($>/\ge$, $</\le$), equality ($=$\,$=$\,/\,$\neq$), boolean ($\wedge/\vee$),
boolean constant inversion, and numeric constant increment --- restricted to the target's
line range and capped per function by language (20 mutants for every Java suite, 30 for every
Python suite). The cap binds two Java functions and no Python one
(Section~\ref{sec:threats}); it does not affect T4, which asks only whether any mutant is
killed.

Two procedural safeguards are worth stating because their absence silently corrupts results.
First, mutation requires overwriting the source; we mutate a \emph{copy} of the source tree
placed ahead of the original on the import path, so that concurrent measurements reading the
same tree are unaffected. Second, tool versions must be matched across families: JUnit~5
pairs Jupiter~5.x with Platform~1.x, and mixing a Jupiter~6 artifact with JUnit~5-style tests
makes the launcher discover zero tests while still exiting successfully --- a failure that
reports as a legitimate zero.

\subsection{Baselines}\label{sec:baselines}

Each baseline generates per class (Java) or per module (Python), matching how the tools
operate; results are then attributed to individual functions by line range, exactly as for
the LLM suites.

\begin{description}
  \item[EvoSuite 1.2.0], 60\,s per class, run under JDK~11. EvoSuite performs deep reflection
    into JDK internals, which the module system blocks from Java~9 onward, and its master
    process launches its client from \texttt{JAVA\_HOME} --- so both must point at the same
    JDK, or the client dies while the master still exits successfully. Because the projects
    were compiled with JDK~17, we recompiled sources to Java~11 bytecode into a separate
    output directory, leaving the original build untouched.
  \item[Randoop 4.3.3], 60\,s per class, JDK~17.
  \item[Pynguin 0.45.0], 90\,s per module, under \textbf{two configurations reported side by
    side}. In its default configuration Pynguin serialises module state to a worker process
    with \texttt{dill} and aborts on C-extension type objects that cannot be resolved by
    name, producing no output at all for several modules. We therefore also run it with the
    master--worker split disabled. Reporting both separates ``the tool scored low'' from
    ``the tool never ran''.
\end{description}

\subsection{Statistical Analysis}\label{sec:stats}

Conditions are compared with the paired Wilcoxon signed-rank test, matching on function
identifier, at $\alpha = 0.05$, with matched-pairs rank-biserial correlation as effect size.
Java and Python are analysed separately throughout: the toolchains differ, and Java enforces
access control while Python does not, so their failure modes are not commensurable. Every
rate is reported over the functions of its own language ($n = 59$ for Java after the
exclusion above, $n = 60$ for Python).

Results are given at two levels in parallel: \emph{all} ($n = 59$ for Java and $n = 60$ for
Python, with unmeasurable cases scored $0$) and \emph{effective} (the subset with a non-zero
value). The \emph{all} level is reported even where the \emph{effective} level is more
favourable.

All configuration decisions above --- criteria, tool versions, budgets, the two-branch
Pynguin design, and the hypotheses --- were recorded in a pre-registration committed before
any baseline was executed.
